% !PRINTED 56
\weberpage{56}
%<<CONT

\begin{equation}
\begin{aligned}
 f'(\alpha_1)&=(\alpha_1-\alpha_2)(\alpha_1-\alpha_3)\cdots(\alpha_1-\alpha_n),\\
 f'(\alpha_2)&=(\alpha_2-\alpha_1)(\alpha_2-\alpha_3)\cdots(\alpha_2-\alpha_n),\\
 &\quad\cdots\\
 f'(\alpha_n)&=(\alpha_n-\alpha_1)(\alpha_n-\alpha_2)\cdots(\alpha_n-\alpha_{n-1}).
\end{aligned}
\tag{12}
\end{equation}

\begin{center}
{\large \S.\ 15.}

\vspace{0.5em}
{\large\gesperrt{Ganze Functionen mehrerer Veränderlichen: Formen.}}
\end{center}

Wir haben bisher vorzugsweise die ganzen rationalen Functionen von einer Veränderlichen betrachtet; wir können uns aber nicht immer darauf beschränken und haben ja auch schon oben Functionen mehrerer Veränderlichen benutzt. Unter einer ganzen rationalen Function $n$ten Grades mehrerer Veränderlichen $F(x,y,z,\ldots)$ verstehen wir eine Summe von Gliedern:
\[
 \sum A_{\alpha,\beta,\gamma,\ldots}x^\alpha y^\beta z^\gamma\cdots,
\]
worin $\alpha,\beta,\gamma,\ldots$ ganzzahlige nicht negative Exponenten sind, deren Summe $\alpha+\beta+\gamma+\cdots$ den Werth $n$ nicht übersteigt, und wenigstens in einem Gliede auch wirklich erreicht. Der Grad wird also bestimmt durch den grössten Werth, den die Summe $\alpha+\beta+\gamma+\cdots$ annimmt.

Wenn die Summe der Exponenten $\alpha+\beta+\gamma+\cdots$ in allen Gliedern denselben Werth hat, so heisst die Function \gesperrt{homogen}.

\gesperrt{Eine fundamentale Eigenschaft der homogenen Functionen $n$ten Grades ist die, dass, wenn alle Variablen mit demselben Factor vervielfältigt werden, der Erfolg derselbe ist, wie wenn die Function mit der $n$ten Potenz vervielfältigt wird}; in Zeichen, wenn $t$ eine beliebige Veränderliche bedeutet:
\begin{equation}
 F(tx,ty,tz,\ldots)=t^nF(x,y,z,\ldots);
\tag{1}
\end{equation}
denn ersetzt man in dem Product $x^\alpha y^\beta z^\gamma\cdots$ die Variablen durch $tx,ty,tz,\ldots$, so erhält es den Factor
\[
 t^{\alpha+\beta+\gamma+\cdots};
\]
hat nun $\alpha+\beta+\gamma+\cdots$ in allen Gliedern denselben Werth $n$, so kann der Factor $t^n$ vor die Summe herausgenommen werden. Hat aber die Summe $\alpha+\beta+\gamma+\cdots$ in den einzelnen Gliedern verschiedene Werthe, so kann ein solcher gemeinschaftlicher Factor nicht herausgenommen werden, wenig-
%CONT>>
